\begin{table}[!htb]
    \centering
    \captionsetup{justification=centering}
    \caption{Processo do erro amostral identificado por estrato geográfico e indicador}
    \resizebox{\textwidth}{!}{
    \begin{tabular}{@{}lccc@{}}
        \toprule
        \textbf{Estrato geográfico} & \textbf{Total de desocupados} & \textbf{Total de ocupados} & \textbf{Taxa de desocupação} \\
        \midrule
        01 -- Belo Horizonte & AR(1) & AR(1) & AR(1) \\
        02 -- Entorno e Colar Metropolitano de BH & MA(1) & AR(1) & MA(1) \\
        03 -- Sul de Minas & AR(1) & AR(1) & AR(1) \\
        04 -- Triângulo Mineiro & AR(1) & AR(1) & AR(1) \\
        05 -- Zona da Mata & AR(1) & AR(1) & AR(1) \\
        06 -- Norte de Minas & AR(1) & AR(1) & MA(1) \\
        07 -- Vale do Rio Doce & AR(1) & AR(1) & AR(1) \\
        08 -- Central & AR(3) & AR(1) & AR(3) \\
        \bottomrule
    \end{tabular}
    }
    \fonte{Elaboração própria, com base nos dados da PNAD Contínua. Nota: a identificação é individual por estrato e indicador. Entre os candidatos cujos resíduos não apresentam autocorrelação significativa pelo teste de Ljung-Box a 5\%, seleciona-se o mais parcimonioso, com desempate pelo BIC.}
    \label{tab:modelos_arma}
\end{table}
